\documentclass[12pt, a4paper]{article}

\usepackage[utf8]{inputenc}
\usepackage[spanish]{babel}
\usepackage{geometry}
\usepackage{titlesec}
\usepackage{hyperref}
\usepackage{enumitem}
\usepackage{xcolor}

\geometry{margin=2.5cm}

\definecolor{cabecera}{RGB}{30, 90, 160}

\titleformat{\section}{\large\bfseries\color{cabecera}}{}{0em}{}[\titlerule]

% --------------------------------------------------
\begin{document}

\begin{center}
    {\LARGE \textbf{Reparto de Tareas}}\\[0.4em]
    {\large \textit{Space Invaders — Implementación de Patrones de Diseño}}\\[0.3em]
    {\small \textit{Equipo: Cachopín}}\\[0.2em]
    {\small Fecha: \today}
\end{center}

\vspace{1em}
\hrule
\vspace{1.5em}

% --------------------------------------------------
\section{Miembros del equipo}

\begin{itemize}[leftmargin=2em]
    \item \textbf{Miembro 1:} Unai Urrutia
    \item \textbf{Miembro 2:} Ekain Ortega
    \item \textbf{Miembro 3:} Eder Lorenzo
    \item \textbf{Miembro 4:} Asier López
\end{itemize}

\vspace{1.5em}

% --------------------------------------------------
\section{Planificación inicial}

\textit{Tareas previstas para la implementación de patrones de diseño y su tiempo estimado:}

\begin{enumerate}[leftmargin=2em]
    \item Implementar patrón Factory para creación de naves — [3] horas
    \item Implementar patrón Strategy para tipos de disparo — [6] horas
    \item Implementar patrón Composite para estructura jerárquica de píxeles — [3] horas
    \item Actualizar diagrama de clases UML — [5] horas
    \item Crear 3 tipos de naves diferentes (Nave1, Nave2, Nave3) — [2] horas
    \item Implementar naves multipixel con geometría variable — [4] horas
    \item Implementar enemigos multipixel con estructura Composite — [1] horas
    \item Implementar tipos de disparo (Pixel, Flecha, Rombo) — [1] horas
    \item Documentar patrones Factory, Strategy y Composite — [2] horas
\end{enumerate}

\textbf{Tiempo total estimado:} [27] horas.

\vspace{1.5em}

% --------------------------------------------------
\section{Resumen de tareas realizadas}

\textit{Tareas efectivamente completadas y su tiempo real:}

\begin{enumerate}[leftmargin=2em]
    \item Implementar patrón Factory para creación de naves — [4] horas
    \item Implementar patrón Strategy para tipos de disparo — [8] horas
    \item Implementar patrón Composite para estructura jerárquica de píxeles — [6] horas
    \item Actualizar diagrama de clases UML — [8] horas
    \item Crear 3 tipos de naves diferentes (Nave1, Nave2, Nave3) — [5] horas
    \item Implementar naves multipixel con geometría variable — [6] horas
    \item Implementar enemigos multipixel con estructura Composite — [3] horas
    \item Implementar tipos de disparo (Pixel, Flecha, Rombo) — [5] horas
    \item Documentar patrones Factory, Strategy y Composite — [5] horas
\end{enumerate}

\textbf{Tiempo total real:} [50] horas.

\vspace{1.5em}

% --------------------------------------------------
\section{Detalle de tareas}

% ---- Tarea 1 ----
\subsection*{Tarea 1 — Implementar patrón Factory para creación de naves}
\begin{description}[leftmargin=2em, style=nextline]
    \item[\textbf{Descripción:}] Desarrollo de la clase \texttt{NaveFactory} como punto único de creación de naves jugables. Implementación del patrón Singleton para garantizar una única instancia. El método \texttt{generate(String)} recibe un identificador y retorna la nave correspondiente con parámetros predefinidos.
    \item[\textbf{Tiempo estimado:}] [3] horas.
    \item[\textbf{Tiempo real:}] [4] horas.
    \item[\textbf{Responsables y dedicación:}]~
        \begin{itemize}[leftmargin=1em, nosep]
            \item Unai Urrutia — [0.5] horas
            \item Ekain Ortega — [2] horas
            \item Eder Lorenzo — 0 horas
            \item Asier López — [1.5] horas
        \end{itemize}
\end{description}

\vspace{0.8em}

% ---- Tarea 2 ----
\subsection*{Tarea 2 — Implementar patrón Strategy para tipos de disparo}
\begin{description}[leftmargin=2em, style=nextline]
    \item[\textbf{Descripción:}] Creación de la interfaz \texttt{StrategyDisparo} con métodos: \texttt{getTipo()}, \texttt{getMunicion()}, \texttt{gastar()}, \texttt{tieneMunicion()} y \texttt{crearDisparo()}. Implementación de tres estrategias concretas: \texttt{DisparoPixel} (munición infinita), \texttt{DisparoFlecha} (30 unidades) y \texttt{DisparoRombo} (20 unidades).
    \item[\textbf{Tiempo estimado:}] [6] horas.
    \item[\textbf{Tiempo real:}] [8] horas.
    \item[\textbf{Responsables y dedicación:}]~
        \begin{itemize}[leftmargin=1em, nosep]
            \item Unai Urrutia — [1] horas
            \item Ekain Ortega — [3] horas
            \item Eder Lorenzo — 0 horas
            \item Asier López — [4] horas
        \end{itemize}
\end{description}

\vspace{0.8em}

% ---- Tarea 3 ----
\subsection*{Tarea 3 — Implementar patrón Composite para estructura jerárquica de píxeles}
\begin{description}[leftmargin=2em, style=nextline]
    \item[\textbf{Descripción:}] Desarrollo de la interfaz \texttt{Component} que define el contrato común. Implementación de \texttt{Pixel} como hoja del Composite y \texttt{Composite} como contenedor de componentes. Diferenciación entre píxeles de nave (proyectil=false) y píxeles de disparo (proyectil=true) con comportamientos distintos. Implementación de métodos de movimiento validado, notificación de eventos y gestión de ciclo de vida.
    \item[\textbf{Tiempo estimado:}] [3] horas.
    \item[\textbf{Tiempo real:}] [6] horas.
    \item[\textbf{Responsables y dedicación:}]~
        \begin{itemize}[leftmargin=1em, nosep]
            \item Unai Urrutia — [2] horas
            \item Ekain Ortega — [2] horas
            \item Eder Lorenzo — 0 horas
            \item Asier López — [2] horas
        \end{itemize}
\end{description}

\vspace{0.8em}

% ---- Tarea 4 ----
\subsection*{Tarea 4 — Actualizar diagrama de clases UML}
\begin{description}[leftmargin=2em, style=nextline]
    \item[\textbf{Descripción:}] Elaboración del diagrama UML de clases que muestre las relaciones entre \texttt{NaveFactory}, \texttt{Naves} (clase abstracta), \texttt{Nave1}, \texttt{Nave2}, \texttt{Nave3}, \texttt{StrategyDisparo} y sus implementaciones, \texttt{Component}, \texttt{Pixel}, \texttt{Composite}, \texttt{Disparo}, \texttt{Espacio} y otros elementos del modelo. Se debe reflejar la jerarquía de herencia, las relaciones de composición y la implementación de interfaces.
    \item[\textbf{Tiempo estimado:}] [5] horas.
    \item[\textbf{Tiempo real:}] [8] horas.
    \item[\textbf{Responsables y dedicación:}]~
        \begin{itemize}[leftmargin=1em, nosep]
            \item Unai Urrutia — [6] horas
            \item Ekain Ortega — [1] horas
            \item Eder Lorenzo — 0 horas
            \item Asier López — [1] horas
        \end{itemize}
\end{description}

\vspace{0.8em}

% ---- Tarea 5 ----
\subsection*{Tarea 5 — Crear 3 tipos de naves diferentes (Nave1, Nave2, Nave3)}
\begin{description}[leftmargin=2em, style=nextline]
    \item[\textbf{Descripción:}] Desarrollo de las clases \texttt{Nave1}, \texttt{Nave2} y \texttt{Nave3} que extienden \texttt{Naves}. Cada nave tiene una combinación diferente de estrategias de disparo permitidas: Nave1 (Pixel y Flecha), Nave2 (Pixel y Rombo), Nave3 (Pixel, Flecha y Rombo). Cada nave tiene geometría única y origen de disparo específico.
    \item[\textbf{Tiempo estimado:}] [2] horas.
    \item[\textbf{Tiempo real:}] [5] horas.
    \item[\textbf{Responsables y dedicación:}]~
        \begin{itemize}[leftmargin=1em, nosep]
            \item Unai Urrutia — [2] horas
            \item Ekain Ortega — [1] horas
            \item Eder Lorenzo — 0 horas
            \item Asier López — [2] horas
        \end{itemize}
\end{description}

\vspace{0.8em}

% ---- Tarea 6 ----
\subsection*{Tarea 6 — Implementar naves multipixel con geometría variable}
\begin{description}[leftmargin=2em, style=nextline]
    \item[\textbf{Descripción:}] Implementación del método \texttt{construir()} en cada nave específica. Cada nave está compuesta por un \texttt{Composite} raíz que contiene múltiples \texttt{Pixel} en posiciones específicas formando formas geométricas distintas. Nave1 forma una T invertida, Nave2 un rectángulo, Nave3 una T normal. Implementación de métodos como \texttt{celdasOcupadas()} para retornar la geometría de cada nave.
    \item[\textbf{Tiempo estimado:}] [4] horas.
    \item[\textbf{Tiempo real:}] [6] horas.
    \item[\textbf{Responsables y dedicación:}]~
        \begin{itemize}[leftmargin=1em, nosep]
            \item Unai Urrutia — [1] horas
            \item Ekain Ortega — [2.5] horas
            \item Eder Lorenzo — 0 horas
            \item Asier López — [2.5] horas
        \end{itemize}
\end{description}

\vspace{0.8em}

% ---- Tarea 7 ----
\subsection*{Tarea 7 — Implementar enemigos multipixel con estructura Composite}
\begin{description}[leftmargin=2em, style=nextline]
    \item[\textbf{Descripción:}] Desarrollo de la clase \texttt{Enemigo} usando el patrón Composite para crear enemigos compuestos por múltiples píxeles. Los enemigos tienen geometría multipixel similar a las naves. Implementación de métodos de movimiento, rotación y comportamiento táctico en flota. Integración con \texttt{FlotaEnemigos} para manejo coordinado de múltiples enemigos.
    \item[\textbf{Tiempo estimado:}] [1] horas.
    \item[\textbf{Tiempo real:}] [3] horas.
    \item[\textbf{Responsables y dedicación:}]~
        \begin{itemize}[leftmargin=1em, nosep]
            \item Unai Urrutia — [0.5] horas
            \item Ekain Ortega — [1] horas
            \item Eder Lorenzo — 0 horas
            \item Asier López — [1,5] horas
        \end{itemize}
\end{description}

\vspace{0.8em}

% ---- Tarea 8 ----
\subsection*{Tarea 8 — Implementar tipos de disparo (Pixel, Flecha, Rombo)}
\begin{description}[leftmargin=2em, style=nextline]
    \item[\textbf{Descripción:}] Implementación del método \texttt{crearDisparo()} en cada estrategia concreta. DisparoPixel crea un Pixel simple (1 celda). DisparoFlecha crea un Composite con 3 píxeles en forma de V invertida. DisparoRombo crea un Composite con 5 píxeles en forma de rombo. Cada disparo es un \texttt{Component} que se comporta como proyectil (notifica movimiento, se desactiva al salir del tablero).
    \item[\textbf{Tiempo estimado:}] [1] horas.
    \item[\textbf{Tiempo real:}] [5] horas.
    \item[\textbf{Responsables y dedicación:}]~
        \begin{itemize}[leftmargin=1em, nosep]
            \item Unai Urrutia — [2] horas
            \item Ekain Ortega — [1] horas
            \item Eder Lorenzo — 0 horas
            \item Asier López — [2] horas
        \end{itemize}
\end{description}

\vspace{0.8em}

% ---- Tarea 9 ----
\subsection*{Tarea 9 — Documentar patrones Factory, Strategy y Composite}
\begin{description}[leftmargin=2em, style=nextline]
    \item[\textbf{Descripción:}] Redacción de documentación completa en formato LaTeX describiendo la implementación de los tres patrones de diseño. Se incluye: descripción general de cada patrón, problema que resuelve, implementación en el proyecto, interfaces y clases concretas, diagramas de secuencia, flujos de ejecución, ventajas, y cómo trabajan juntos los patrones. Elaboración de diagramas conceptuales para ilustrar las relaciones entre componentes.
    \item[\textbf{Tiempo estimado:}] [2] horas.
    \item[\textbf{Tiempo real:}] [5] horas.
    \item[\textbf{Responsables y dedicación:}]~
        \begin{itemize}[leftmargin=1em, nosep]
            \item Unai Urrutia — [1] horas
            \item Ekain Ortega — [3] horas
            \item Eder Lorenzo — 0 horas
            \item Asier López — [1] horas
        \end{itemize}
\end{description}

\vspace{0.8em}

% --------------------------------------------------
\section{Notas y observaciones}

\begin{itemize}[leftmargin=2em]
    \item \textbf{Desviaciones significativas:} El tiempo real (50 horas) superó en 85\% la estimación inicial (27 horas). Las tareas de mayor complejidad como Strategy y Composite requirieron más tiempo del previsto debido a la integración entre patrones y la necesidad de refactorizar código existente.
    \item \textbf{Problemas encontrados:} Dificultades en la coordinación entre patrones Composite y Strategy para disparos complejos. Problemas de rendimiento con notificaciones múltiples del Observer. Complejidad inesperada en la validación de movimientos de naves multipixel. Tiempo adicional requerido para actualizar diagramas UML debido a cambios iterativos en el diseño.
    \item \textbf{Lecciones aprendidas:} La importancia de diseñar interfaces claras antes de implementar patrones concretos. Los patrones de diseño requieren más tiempo de planificación inicial pero facilitan enormemente las extensiones posteriores. La documentación simultánea al desarrollo es más eficiente que documentar al final.
    \item \textbf{Sugerencias para futuros proyectos:} Incluir más tiempo de buffer para integración entre patrones (al menos 40\% adicional). Realizar pruebas incrementales de cada patrón antes de integrar. Considerar el uso de herramientas de diagramado automático para mantener UML actualizado.
\end{itemize}

\vspace{1.5em}

% --------------------------------------------------
\section{Resumen estadístico}

\begin{center}
    \begin{tabular}{|l|r|r|}
        \hline
        \textbf{Miembro} & \textbf{Horas estimadas} & \textbf{Horas reales} \\
        \hline
        Unai Urrutia & [8.5] & [16] \\
        \hline
        Ekain Ortega & [9.2] & [16.5] \\
        \hline
        Eder Lorenzo & 0 & 0 \\
        \hline
        Asier López & [9.5] & [17.5] \\
        \hline
        \textbf{TOTAL} & \textbf{[27.2]} & \textbf{[50]} \\
        \hline
    \end{tabular}
\end{center}

\end{document}
